\documentclass[11pt]{article}

\usepackage[a4paper,margin=2.5cm]{geometry}
\usepackage[T1]{fontenc}
\usepackage[utf8]{inputenc}
\usepackage{lmodern}
\usepackage{microtype}
\usepackage{xcolor}
\usepackage{amsmath,amssymb}
\usepackage{enumitem}
\usepackage[hidelinks]{hyperref}

\setlength{\parindent}{0pt}
\setlength{\parskip}{0.7em}
\setlist[itemize]{leftmargin=*,itemsep=0.2em,topsep=0.3em}

\definecolor{reviewgray}{RGB}{70,70,70}
\definecolor{responseblue}{RGB}{0,72,125}

\newenvironment{reviewercomment}{\begin{quote}\color{reviewgray}\itshape}{\end{quote}}
\newenvironment{authorresponse}{\par\noindent\textbf{\color{responseblue}Response.}\ }{\par}
\newcommand{\change}[1]{\par\noindent\textit{Changes in the revised manuscript.} #1\par}

\begin{document}

\begin{center}
{\Large\bfseries Response to the Editor and Reviewers}\\[0.5em]
{\itshape Manuscript: ``An attractor account of numerosity perception''}\\[0.25em]
{\itshape Manuscript ID: [insert manuscript ID]}
\end{center}

Dear Editor and Reviewers,

We thank the reviewers for their careful, constructive, and technically detailed assessments. Their comments led us to substantially rework both the model and the empirical analysis. The principal revisions are as follows:

\begin{itemize}
\item We replaced the summary-curve objective with a trial-level categorical likelihood over all 52,995 individual integer reports. Response time is now modelled as a separate trial-level log-normal channel.
\item We replaced the original Poisson-well parameterisation with integer-centred Gaussian memory components whose widths increase with magnitude according to a power law.
\item The primary model now uses a single attractor landscape shared across all five exposure durations; only the sensory channel and lapse vary by duration.
\item We added matched static, constant-width, and purely driven-accumulator controls to isolate the contributions of magnitude-dependent broadening, attractor clean-up, and online accumulation.
\item Response time is now derived from bounded attractor-convergence time. Residual entropy was retained only as a comparison and did not improve the selected shared RT mapping.
\item We removed the proposed density regime and its comparison with density-threshold studies, and we restricted the central claim to a shared numerical readout spanning more categorical and more approximate responses.
\item We expanded the discussion of competing visuospatial-capacity, recurrent visual, image-computable, and other dynamical models.
\item We added a bivariate hierarchical individual-differences analysis using non-overlapping lower- and higher-numerosity trials. Participants with greater fitted landscape resolution also show lower Weber noise.
\item We reformulated and strengthened the exploratory developmental-dyscalculia analysis while making its limitations and extrapolative status explicit.
\end{itemize}

We address each comment below.

\section*{Reviewer 1}

\subsection*{Major Comment 1: The model comparison was based on the wrong likelihood}

\begin{reviewercomment}
The original AIC/BIC comparison was computed from Gaussian residuals on three summary curves. The reviewer requested a trial-level response likelihood for the discrete integer reports, with RT modelled separately.
\end{reviewercomment}

\begin{authorresponse}
We agree. The original comparison discarded important information in the discrete response distributions and did not provide an appropriate basis for the principal model comparison.

We therefore rebuilt the analysis around a categorical trial-level likelihood. For each presented numerosity, the probability of every integer response is obtained by marginalising the model readout over the sensory distribution,
\[
p(r\mid n)=\int p\!\left(r\mid x^*(y_0)\right)p(y_0\mid n)\,dy_0.
\]
The integral is evaluated deterministically, and the likelihood is computed over all 52,995 individual integer reports. No trials are aggregated. Typed responses outside the integer-well support are retained through the same uniform lapse channel in every model, and the lapse is included in all AIC and BIC calculations.

Response times are no longer combined with response summaries in a single loss. They are modelled separately using a trial-level log-normal likelihood, with the landscape fixed at the values obtained from the categorical-response fit. The RT analysis is therefore a constrained cross-channel test of the response-fitted mechanism.
\end{authorresponse}

\change{The revised Methods and Supplementary Material 1 now specify the categorical response likelihood, deterministic quadrature, common lapse convention, well and response supports, and separate RT channel. The resulting model comparisons are reported in Figure~2a and Supplementary Table~S2.}

\subsection*{Major Comment 2: The reported AIC differences were too small}

\begin{reviewercomment}
The original summary-level difference between the Hopfield--Poisson model and its no-clean-up control was only $\Delta\mathrm{AIC}=1.6$, which did not support strong claims about the contribution of the dynamics.
\end{reviewercomment}

\begin{authorresponse}
We agree that the original summary-level differences did not support strong conclusions about model preference. This was one of the principal motivations for replacing the original analysis.

Under the revised trial-level likelihood, the model-comparison margins are substantially larger. Relative to the power-law attractor, the matched static power-law readout trails by $\Delta\mathrm{AIC}=722$, the constant-width attractor by $\Delta\mathrm{AIC}=3488$, the Weber/ANS model by $\Delta\mathrm{AIC}=838$, the matched bounded-optimal model by $\Delta\mathrm{AIC}=11{,}453$, and the purely driven accumulator by $\Delta\mathrm{AIC}=7{,}954$. The corresponding BIC comparison gives the same model ordering.

Importantly, the static control has one fewer parameter because it does not require the shared stopping budget $G$. The attractor advantage therefore remains after penalising this additional flexibility.
\end{authorresponse}

\change{The manuscript now bases its empirical case on the trial-level comparison in Figure~2a and Supplementary Table~S2 rather than on the original summary-curve AIC values.}

\subsection*{Major Comment 3: The Hopfield framing was largely decorative}

\begin{reviewercomment}
The reviewer questioned what the one-dimensional attractor dynamics contributed beyond a static mixture or noisy MAP readout, and noted that the model did not address the high-dimensional storage-capacity properties usually associated with modern Hopfield networks.
\end{reviewercomment}

\begin{authorresponse}
We agree that the original analysis did not adequately isolate what the attractor dynamics contributed beyond the static response family.

The revised manuscript contains a matched static control that retains the same sensory channel, integer-centred components, width law, and categorical readout but removes gradient-descent clean-up. Despite requiring one fewer parameter, this static model fits substantially worse than the attractor model. We also include constant-width versions of the static and attractor models, which separately test whether the advantage comes from the dynamics or simply from magnitude-dependent broadening.

The dynamics now make a second distinctive prediction through convergence time. The selected RT model maps the expected number of bounded clean-up steps onto response time; an entropy-only model performs substantially worse. Thus, the chronometric analysis now depends directly on the operation that differentiates the attractor from its static twin.

We have also narrowed the theoretical use of the Hopfield framework. We do not claim the high-dimensional storage capacity associated with modern Hopfield networks. The relevant claim is a one-dimensional associative-retrieval operation in which a noisy scalar state is iteratively transformed by a log-sum-exp energy landscape generated by stored integer-centred components. The revised manuscript therefore describes the model as Hopfield-inspired and focuses on pattern completion rather than capacity.
\end{authorresponse}

\change{The matched static and constant-width controls are shown in Figure~2a and Supplementary Table~S2; the convergence-time analysis is shown in Figures~1c and 2b.}

\subsection*{Major Comment 4: The accumulator interpretation of the Poisson form was not derived}

\begin{reviewercomment}
The reviewer argued that the original Poisson wells were a parametric choice rather than a genuine implementation of a classical accumulator, and that the manuscript used the accumulator label more strongly than the formal model justified.
\end{reviewercomment}

\begin{authorresponse}
We agree. We have removed the original Poisson-well parameterisation and no longer describe the main model as an accumulator model.

The revised landscape is generated by Gaussian memory components centred on the integers, with widths
\[
\sigma_n=\sigma_c n^p.
\]
The power-law broadening is motivated by the decline in number-word frequency, rather than being presented as a derivation from accumulation statistics. The model now begins with an already formed scalar numerosity estimate $y_0$, and the attractor implements subsequent clean-up and numerical readout.

We also tested the stronger version of the original accumulator intuition explicitly. In a purely driven model, registered items generate a running count $K(t)$ that drives the attractor during exposure. This model performs substantially worse than the completed-estimate attractor. The result clarifies the division of labour: a partial count is itself a valid intermediate cardinality and is not diagnostic of the final total, so the attractor is repeatedly drawn towards intermediate integer memories. We now interpret the attractor as a mechanism operating on the output of accumulation rather than as the accumulation process itself.
\end{authorresponse}

\change{The Poisson formulation and accumulator claims have been removed. The purely driven running-count model is now a quantitative control described in Methods, Results, Discussion, and Supplementary Figure~S1.}

\subsection*{Major Comment 5: RT was not derived from the attractor dynamics}

\begin{reviewercomment}
The original RT model used posterior entropy. The reviewer suggested deriving RT from convergence time or first-passage time on the attractor landscape.
\end{reviewercomment}

\begin{authorresponse}
We agree and have replaced the original entropy-based RT account.

For each sensory input, we now record the number of clean-up iterations $T$. Because trajectories in the flat large-number region may not converge, the process is bounded by a global stopping budget $G$,
\[
T=\min(T^*,G).
\]
The budget is selected by profiling the categorical-response likelihood, is shared across durations, and is counted as one free parameter. The selected value is $G=120$.

The selected response-time model is
\[
RT(n)=t_0+\tau\,\mathbb{E}[T(n)],
\]
with one $\tau$ shared across exposure durations. The landscape is held fixed at its categorical-response fit before the RT channel is estimated.

We compared this model with an entropy-only model and a hybrid convergence-plus-entropy model. The convergence-only and hybrid models were effectively indistinguishable ($\Delta\mathrm{AIC}=0.24$), so we retained convergence alone as the more parsimonious account. The entropy-only model was much worse ($\Delta\mathrm{AIC}=1579.5$). Thus, under a common mapping across durations, static posterior uncertainty adds no supported predictive value beyond the attractor's convergence time.
\end{authorresponse}

\change{The exact integration rule, convergence and oscillation criteria, step clamp, stopping budget, empirical RT origin, log-normal likelihood, and fitted RT parameters are now documented in Methods and Supplementary Material~1.}

\subsection*{Major Comment 6: The three-regime and subitizing claims were overstated}

\begin{reviewercomment}
The reviewer noted that the density regime was not modelled and argued that fitting small-number accuracy did not amount to a mechanistic account of all classical subitizing phenomena, because the model contains no spatial, object-based, or attentional mechanism.
\end{reviewercomment}

\begin{authorresponse}
We agree that the original manuscript conflated three distinct claims: fitting small-number behaviour, providing a complete mechanism of subitizing, and predicting a transition to density perception.

We have removed the density regime entirely, including the proposed density transition and the comparison with Anobile et al. The revised paper concerns the change from more categorical to increasingly approximate numerical reports within the fitted range.

We have also narrowed the subitizing claim. Under the revised trial-level likelihood, the model reproduces near-ceiling exact-response rates for the smallest numerosities, but we agree that this does not constitute a model of all classical subitizing phenomena. The attractor has no representation of object locations, configurations, or attentional load. We therefore describe it as a model of the numerical readout downstream of modality-specific sensory and attentional processing, and explicitly state that spatial individuation and attention may make additional contributions to subitizing.

We retain a narrower unification claim: the same numerical readout produces more categorical reports where integer components are well separated and Weber-like reports where they overlap. This does not require every perceptual signature of subitizing to arise from the landscape, nor does it rule out object-individuation processes for the smallest sets.
\end{authorresponse}

\change{All density-regime analyses and claims have been removed. The Abstract, Introduction, Discussion, and Conclusion now restrict the proposal to a shared numerical readout, and the Discussion explicitly identifies spatial and attentional subitizing effects as outside the present model.}

\subsection*{Major Comment 7: Insufficient engagement with competing literature}

\subsubsection*{Cheyette, Wu, and Piantadosi (2024)}

\begin{reviewercomment}
The reviewer asked us to engage directly with the visuospatial-capacity account, which attributes number psychophysics to bandwidth-limited visual memory and reports within-participant transfer between non-numerical visual memory and numerosity estimation.
\end{reviewercomment}

\begin{authorresponse}
We agree that the original treatment was insufficient. The revised manuscript now presents the attractor and visuospatial-capacity accounts as competing explanations of where the same psychophysical regularities arise.

The visuospatial account places the principal limitation in sensory encoding and predicts shared individual constraints across non-numerical visual-memory and numerosity tasks. We now discuss the within-participant transfer of information-capacity estimates from change localisation to number estimation reported by Cheyette et al. (2024).

The attractor account instead predicts that a shared numerical landscape should generalise across sensory modalities and may be selectively modified by numerical experience. We formulate differential tests involving cross-modal parameter sharing, numerical versus visuospatial training, and developmental differences. We also discuss evidence that nonsymbolic acuity can remain robust with little formal schooling, which constrains a strongly schooling-dependent interpretation of the landscape.

A new bivariate hierarchical analysis provides an additional test within the present account. Using non-overlapping trials, landscape resolution estimated from lower numerosities predicts Weber noise at larger numerosities after controlling for exposure duration and finite-trial uncertainty ($\hat\rho=-.569$, $N=362$).
\end{authorresponse}

\change{The competing explanatory loci and their differential predictions are now discussed in the Introduction and Discussion; the individual-differences analysis is reported in Figure~2c and Supplementary Material~3.}

\subsubsection*{Recurrent visual-network literature}

\begin{reviewercomment}
The reviewer asked us to engage with the recurrent on-centre/off-surround models developed by Sengupta and colleagues, including their range effects, adaptation, and switch-cost predictions.
\end{reviewercomment}

\begin{authorresponse}
We have expanded the Introduction and Discussion to acknowledge the recurrent visual-network literature, including Sengupta et al. (2014), Sengupta et al. (2017), and Verma and Sengupta (2023), and we no longer claim novelty for a three-range or density boundary.

We now distinguish these accounts from the present proposal by explanatory level. The recurrent visual models derive numerical sensitivity and range effects from the dynamics of a sensory network, whereas the present model begins with scalar numerical evidence and asks how that evidence is stabilised by an explicit memory landscape. We also identify adaptation and switch-cost phenomena as important tests not addressed by the present dataset.
\end{authorresponse}

\change{The relevant recurrent-network literature and the distinction between sensory-network dynamics and central numerical readout have been added to the literature review and Discussion.}

\subsection*{Major Comment 8: Independent fits by duration under-tested model stability}

\begin{reviewercomment}
The reviewer suggested a constrained model in which the representational landscape is shared across exposure durations and only sensory parameters vary.
\end{reviewercomment}

\begin{authorresponse}
We agree. The primary analysis is now a constrained joint model with a single power-law landscape shared across all five exposure durations. Specifically, $\sigma_c$, $p$, and the global stopping budget $G$ are shared, while the sensory-noise parameters and lapse vary with exposure duration.

This is now the headline model throughout the manuscript. It provides the best trial-level account despite being more constrained than the fully duration-specific Weber/ANS model. The shared-landscape attractor beats the Weber/ANS model by $\Delta\mathrm{AIC}=838$ and the matched bounded-optimal model by $\Delta\mathrm{AIC}=11{,}453$.

We therefore no longer infer that the cognitive system uses a different numerical landscape at each duration. Exposure duration is represented through the sensory channel, while the numerical landscape is treated as a stable representational property.
\end{authorresponse}

\change{The sharing structure and parameter counts are specified in Methods, Supplementary Material~1, and Supplementary Table~S2.}

\subsection*{Major Comment 9: The developmental-dyscalculia extension was too weak}

\begin{reviewercomment}
The reviewer argued that the original group-level analysis and its interpretation were too strong for the small dataset and limited numerosity range, and questioned the exclusion of $n=8$.
\end{reviewercomment}

\begin{authorresponse}
We agree that the original group-level fit could not support the causal interpretation we gave it. We have substantially revised the analysis and moderated its conclusions.

First, we rechecked the original study and corrected the task description: the arrays were shown for a fixed 500-ms exposure. The exclusion of $n=8$ follows the original authors' treatment of the guessing end-effect and is now stated explicitly; it was not introduced post hoc for the present model.

Second, rather than fitting only one model per group, we fitted two restricted single-mechanism variants to every participant's accuracy curve: a width-only model, in which broader and more overlapping memory components carry the decline, and a noise-only model, in which the width profile is fixed and sensory noise carries the decline. The width-only variant fits better in all 64 participants (mean SSE $=.010$ versus $.097$; Wilcoxon $p=3.5\times10^{-12}$).

Third, participant-level fitted width profiles are coarser in the developmental-dyscalculia group. The extrapolated width index is lower in DD than in age-matched controls (median $8.1$ versus $8.8$; Mann--Whitney $p=.018$).

We nevertheless retain this section explicitly as exploratory. We now state that the data contain only seven trials per cell, that the fitted width indices extend beyond the observed $n=1$--$7$ range and are therefore extrapolated summaries, that the restricted comparison does not fully identify sensory and representational mechanisms, and that a sensory-noise contribution may coexist. We no longer claim that the analysis explains the intact-subitizing/impaired-estimation contrast or definitively localises the disorder.
\end{authorresponse}

\change{The revised analysis and its limitations are reported in Supplementary Material~4 and Supplementary Figure~S4.}

\subsection*{Minor Comment: Benford framing}

\begin{reviewercomment}
Dehaene and Mehler (1992) report a power-law decline in number-word frequency, not strict Benford behaviour.
\end{reviewercomment}

\begin{authorresponse}
We agree. All references to Benford's law have been removed. The revised model uses a simple power-law width schedule motivated by the cross-linguistic power-law decline in number-word frequency. We do not claim that the empirical frequencies follow strict Benford behaviour.
\end{authorresponse}

\subsection*{Minor Comment: The 40-ms condition}

\begin{reviewercomment}
The reviewer requested a more candid account of what fails at 40 ms rather than describing the condition only as noisier.
\end{reviewercomment}

\begin{authorresponse}
We agree. The revised manuscript now identifies the 40-ms condition as the principal remaining difficulty. Both the mean-estimate and response-time fits are weaker at 40 ms than at the longer exposures, and the fitted sensory channel is markedly noisier. We report these as systematic residuals rather than treating the condition merely as noisy data.
\end{authorresponse}

\subsection*{Minor Comment: Comparison with Anobile et al.}

\begin{reviewercomment}
The reviewer regarded the comparison between exposure-dependent model thresholds and patch-size-dependent density thresholds as numerological.
\end{reviewercomment}

\begin{authorresponse}
We agree. The proposed density transition and the comparison with Anobile et al. have been removed entirely. The revised manuscript does not model density perception.
\end{authorresponse}

\subsection*{Minor Comment: Compression and logarithmic geometry}

\begin{reviewercomment}
The reviewer noted that the original model did not adjudicate between increasing well width on a linear axis and fixed-width wells on a logarithmically compressed axis.
\end{reviewercomment}

\begin{authorresponse}
We agree. The Poisson accumulator parameterisation and the corresponding claim have been removed. The revised paper tests the simplest power-law width schedule motivated by number-word frequency and explicitly states that alternative smooth width functions may provide different descriptions. The central claim is now about magnitude-dependent component broadening and associative clean-up, not about uniquely establishing the origin of number-line compression.
\end{authorresponse}

\section*{Reviewer 2}

\subsection*{Comment 1: Malo et al. (2023)}

\begin{reviewercomment}
The reviewer asked us to acknowledge the quantum-spin model of Malo et al. as a recent Hopfield-like or dynamical account of numerosity.
\end{reviewercomment}

\begin{authorresponse}
Thank you for pointing us to this work. We now cite Malo et al. (2023) in the Introduction and describe it as a distinct dynamical route to a unified numerical representation. We also clarify the difference between the approaches: Malo et al. encode numerosity in the collective spectrum of a quantum-spin model, whereas the present account represents integers as explicit memory components in an attractor landscape and tests their consequences for trial-level categorical reports and response time.
\end{authorresponse}

\subsection*{Comment 2: Modest AIC differences and the source of the improvement}

\begin{reviewercomment}
The reviewer questioned whether the original small advantage reflected the attractor dynamics or simply the Poisson response family.
\end{reviewercomment}

\begin{authorresponse}
We agree that the original summary-level differences could not distinguish the contribution of the Poisson family from that of the dynamics.

Both aspects of the analysis have changed. The Poisson family is no longer used, and all comparisons now use the categorical trial-level likelihood. The contribution of the dynamics is isolated through a matched static readout that retains the same landscape, sensory channel, and response rule but removes gradient-descent clean-up. The attractor improves AIC by 722 despite requiring the additional shared stopping-budget parameter.

The constant-width controls separately show that magnitude-dependent broadening is also necessary: fixing $p=0$ worsens AIC by 3488. The revised comparison therefore distinguishes the two central ingredients of the model: a width profile that broadens with magnitude and attractor clean-up acting on that profile. The RT analysis provides a further test because convergence time, rather than posterior entropy alone, predicts the observed chronometry.
\end{authorresponse}

\subsection*{Comment 3: Parameter identifiability}

\begin{reviewercomment}
The reviewer asked us to highlight the parameter trade-offs more clearly and to avoid strong causal interpretation of individual fitted parameters, especially in the dyscalculia extension.
\end{reviewercomment}

\begin{authorresponse}
We agree. The revised model has a different parameterisation, so the original $c$--$\kappa$ trade-off no longer applies: the Poisson scale $c$ has been removed, and the selected RT model retains convergence time alone ($\kappa=0$).

In the revised landscape, $\sigma_c$ and $p$ jointly define the width profile. We do not interpret them separately, but summarise their combined effect using
\[
n_{\mathrm{width}}=\sigma_c^{-1/p}.
\]
Likewise, the sensory parameters $s_0$ and $\lambda_0$ are interpreted jointly through the resulting sensory-noise function $\sigma(n)$ rather than as independent psychological quantities.

The developmental-dyscalculia analysis has been reformulated accordingly. We refer to a coarser fitted width profile, not to a causal impairment of one independently identified parameter, and the result is explicitly described as exploratory.
\end{authorresponse}

\change{The revised Methods, Results, Supplementary Table~S1, and dyscalculia section consistently interpret width and sensory parameters through their combined profiles rather than as isolated psychological quantities.}

\subsection*{Comment 4: Why fit parameters separately across exposure durations?}

\begin{reviewercomment}
The reviewer asked whether separate fits implied that the numerical mechanism itself changed with exposure duration.
\end{reviewercomment}

\begin{authorresponse}
We agree that the original strategy left this ambiguous. The revised primary model no longer fits a separate landscape at each duration.

A single landscape, defined by shared $\sigma_c$ and $p$, is fitted jointly across all five durations. The global stopping budget $G$ is also shared. Only the sensory channel and lapse vary with duration. This directly tests the hypothesis that the numerical representation is stable while the quality of the available sensory evidence changes with exposure.

The shared-landscape attractor remains the best model in the trial-level comparison, including against a fully duration-specific Weber/ANS code.
\end{authorresponse}

\subsection*{Comment 5: Absence of a perceptual module}

\begin{reviewercomment}
The reviewer asked us to discuss more fully that the model lacks a perceptual front end and therefore cannot simulate effects of continuous image magnitudes, and to engage with image-computable models such as Testolin et al. (2020) and Dolfi et al. (2024).
\end{reviewercomment}

\begin{authorresponse}
We agree and have expanded this limitation substantially.

The revised manuscript defines the attractor as a candidate central numerical-readout stage. The input $y_0$ summarises evidence that has already been extracted by modality-specific processing. Consequently, the model does not simulate dot area, density, convex hull, object segmentation, spatial configuration, or attentional manipulations.

We now discuss image-computable models explicitly, including Stoianov and Zorzi, Testolin et al., and Dolfi et al. A complete perceptual account could use such a model to generate the scalar evidence entering the attractor. This coupled architecture would provide a stronger test of whether attractor clean-up explains response structure beyond that already generated by sensory encoding.

We also formulate cross-modal tests: visual, auditory, and tactile channels could have different sensory parameters while sharing the same numerical landscape. We present this as a prediction rather than a result established by the current visual dataset.
\end{authorresponse}

\change{The revised Introduction and Discussion now distinguish modality-specific sensory extraction from the proposed shared numerical readout, engage with image-computable alternatives, and state explicitly that continuous-magnitude and attentional effects are outside the present model.}

\section*{Closing statement}

We again thank both reviewers for their detailed and constructive comments. The revisions have changed the manuscript substantially: the principal evidence now comes from a trial-level shared-landscape model; the dynamics are isolated by matched controls and linked directly to convergence time; the accumulator and density claims have been removed or tested explicitly; and the theoretical claims have been restricted to the level supported by the data. We believe that the resulting manuscript provides a clearer and more falsifiable account of attractor-based numerical readout.

Sincerely,

Eudald Correig-Fraga

\end{document}
